\documentclass[11pt,a4paper]{article}

\usepackage[utf8]{inputenc}
\usepackage[T1]{fontenc}
\usepackage{lmodern}
\usepackage{geometry}
\usepackage{xcolor}
\usepackage{enumitem}
\usepackage{listings}
\usepackage{graphicx}
\usepackage{parskip}
\usepackage{fancyhdr}
\usepackage[most]{tcolorbox}
\usepackage{titlesec}

% Palette: Dunkelblau and Blau only
\definecolor{Blau}{HTML}{3D5A9B}
\definecolor{Dunkelblau}{HTML}{03003C}

\usepackage[colorlinks=true, linkcolor=Blau, urlcolor=Blau]{hyperref}

\geometry{margin=1in, top=1.2in, bottom=1.2in}

% Header / footer — Dunkelblau background, white text
\pagestyle{fancy}
\fancyhf{}
\lhead{\small\color{Dunkelblau}ML4Health Lecture}
\rhead{\small\color{Dunkelblau}GitHub Classroom Guide}
\cfoot{\small\thepage}
% \renewcommand{\headrulewidth}{0.4pt}
% \renewcommand{\headrule}{\hbox to\headwidth{\color{Blau}\leaders\hrule height\headrulewidth\hfill}}

% Section styling
\titleformat{\section}{\large\bfseries\color{Dunkelblau}}{{\thesection}}{0.8em}{}[\vspace{-0.3em}{\color{Blau}\rule{\linewidth}{0.4pt}}]
\titleformat{\subsection}{\normalsize\bfseries\color{Blau}}{{\thesubsection}}{0.6em}{}

% Shell / code listings
\lstdefinestyle{shell}{
    basicstyle=\ttfamily\small,
    breaklines=true,
    frame=single,
    rulecolor=\color{Blau!40},
    backgroundcolor=\color{Dunkelblau!5},
    xleftmargin=0.5em,
    xrightmargin=0.5em,
    framexleftmargin=0.5em,
    framexrightmargin=0.5em,
    keywordstyle=\color{Blau}\bfseries,
    morekeywords={git, conda, pip, pytest},
    commentstyle=\color{Blau!60},
    showstringspaces=false,
    aboveskip=0.6em,
    belowskip=0.6em,
}
\lstset{style=shell}

% Callout boxes
\tcbset{
    importantbox/.style={
        colback=Dunkelblau!8,
        colframe=Dunkelblau,
        coltitle=white,
        fonttitle=\bfseries,
        boxrule=0.8pt,
        arc=3pt,
        left=6pt, right=6pt, top=4pt, bottom=4pt,
    },
    tipbox/.style={
        colback=Blau!10,
        colframe=Blau,
        coltitle=white,
        fonttitle=\bfseries,
        boxrule=0.8pt,
        arc=3pt,
        left=6pt, right=6pt, top=4pt, bottom=4pt,
    }
}

\title{\textbf{How-To: Work on Assignments in GitHub Classroom}}
\author{ML4Health Lecture Team}
\date{}

\begin{document}

\maketitle
\thispagestyle{fancy}

\section*{About}

This guide explains what is required and how to work on the assignments for the ML4Health lecture. For the installation and set-up of the required libraries, please follow the provided links and consult the respective documentation. As installation is usually platform-specific, we have only limited capacity to help with setup issues. If you are facing issues \emph{accessing} the assignments, please get in touch via the ILIAS forum.

\section{Prerequisites}

You will need the following tools and accounts:

\begin{itemize}[leftmargin=1.5em, itemsep=0.4em]
    \item \href{https://git-scm.com/}{\textbf{git}} — likely pre-installed; verify with \texttt{git --version}. Required to access and submit exercises.
    \item \href{https://github.com}{\textbf{GitHub account}} — required to access assignments via GitHub Classroom.
    \item \href{https://docs.anaconda.com/miniconda/}{\textbf{miniconda}} — verify with \texttt{conda --version}. We recommend miniconda to manage Python environments so that installations for different lectures do not interfere with each other.
\end{itemize}

\section{Accessing and Working on Your Assignment}

Check the ILIAS main page each week for the assignment link. When you enter the classroom for the first time, you must link your GitHub account to one of the provided university identifiers — please select your \textbf{myAccount user ID}.

Upon accepting the assignment, a private repository named

\begin{center}
    \texttt{<assignment-name>-<your-github-name>}
\end{center}

\noindent is created for you.

\subsection{Repository Access}

\begin{tcolorbox}[importantbox, title={Important: Email Invitation Required}]
After accepting the GitHub Classroom invitation link, you may temporarily see a message saying you do not have access to the repository (e.g., ``Repository access issue''). \textbf{This is expected.}

GitHub will send a separate email invitation to your repository or organization. You \emph{must} accept this invitation before you can access your assignment.

If you do not see the email, check your spam folder and confirm you are using the correct address associated with your GitHub account.
\end{tcolorbox}

\begin{center}
    \includegraphics[width=0.75\textwidth]{repository_access_issue.png}
\end{center}

\subsection{Submission and Automated Tests}

Each assignment includes test cases so you can verify your solutions locally:

\begin{lstlisting}
pytest tests/ -v
\end{lstlisting}
\begin{tcolorbox}[tipbox, title={Important: Final Submission}]
Pushing to \texttt{main} triggers the automated test workflow on GitHub. You should treat this as your \textbf{final submission attempt}, so please \textbf{only push to \texttt{main} when you are confident in your solution}.

To save intermediate work remotely, use a separate branch:
\begin{lstlisting}[aboveskip=0.3em, belowskip=0em]
git checkout -b development
\end{lstlisting}
\end{tcolorbox}

\subsection{Feedback and Reflection}

For each assignment, a pull request against a feedback branch will be created automatically. This is how we provide inline feedback on your code. \textbf{Please do not merge this pull request.}

If you were unable to solve an exercise — or solved it but still have questions — please describe the specific point where you got stuck in a file called \texttt{experiences.md} in your repository. General questions should go to the ILIAS forum.

%break to next page
\newpage

\section{Managing Packages with miniconda}

Create a dedicated conda environment for this lecture and activate it:

\begin{lstlisting}
conda create -n ml4health python=3.11
conda activate ml4health
\end{lstlisting}

Each assignment provides its dependencies in a \texttt{requirements.txt}. Install them with:

\begin{lstlisting}
pip install -r requirements.txt
\end{lstlisting}

\section{Access to GPU Compute (bwUniCluster)}

Some assignments — in particular those involving deep neural networks — benefit significantly from GPU acceleration. As a student at a Baden-Württemberg university, you have free access to \textbf{bwUniCluster}, a shared HPC cluster with GPU nodes.

\begin{tcolorbox}[importantbox, title={Register Early}]
Account provisioning can take some time. Please register \emph{before} the first GPU-heavy assignment is released so that you are not blocked when you need compute.
\end{tcolorbox}

\subsection*{How to register}

Follow the official registration guide at:

\begin{center}
    \href{https://wiki.bwhpc.de/e/Registration/bwUniCluster}{\texttt{https://wiki.bwhpc.de/e/Registration/bwUniCluster}}
\end{center}

You will need your university credentials. Once registered, you can log in via SSH and submit jobs through the SLURM scheduler. Further documentation on using the cluster (modules, job scripts, storage) is available on the same wiki.

\end{document}
